\begin{proposition}[The terminal-range form of the original criterion]
\label{prop:v1-terminal}
Fix \(0<\eta<1\). Suppose that, for every \(\delta>0\),
\begin{gather*}
 \E_a(N,H)\ll_{\eta,\delta}N^{2+\delta}/H
 \quad(H\le N^{1-\eta},\ H\text{ dyadic}),\\
 \sum_{N^{1-\eta}<h<N}C^*_{a,h}(N)
        \ll_{\eta,\delta}N^{1+\delta}.
\end{gather*}
Then \(\sup_\alpha|A_N(\alpha)|\ll_{\eta,\varepsilon}
N^{1/2+\varepsilon}\) for every \(\varepsilon>0\).
In particular, for \(a=\lambda\) this implies RH.
\end{proposition}
\begin{proof}
The term \(h=1\) is at most \(N\). In a dyadic block beginning
at \(H\le N^{1-\eta}\), Cauchy--Schwarz gives
\[
 \sum_{H<h\le2H}C^*_{a,h}(N)
 \le\sqrt H\,\E_a(N,H)^{1/2}
 \ll_{\eta,\delta}N^{1+\delta/2}.
\]
There are \(O(\log(2N))\) such blocks. Include the entire last
block, even if its upper end exceeds \(N^{1-\eta}\). All shifts
not covered by these blocks or by \(h=1\) exceed that threshold,
so their sum is bounded by the assumed terminal sum. Equation
\eqref{eq:square} therefore bounds the squared exponential sum by
\(O_{\eta,\delta}(N^{1+\delta/2}\log(2N)+N^{1+\delta})\).
Taking \(0<\delta<\varepsilon\), enlarging the constant for bounded
\(N\), and taking square roots proves the claim. The implication
for \(\lambda\) follows from Theorem~\ref{thm:equiv}.
\end{proof}

The proposition retains the terminal-range formulation of
version~1. Its hypotheses concern the signed Liouville sequence
when that sequence is substituted; they are not established here.
Theorem~\ref{thm:grh} below gives the stronger consequence of
assuming the collar bound over the entire dyadic range.
